% !PRINTED 618
\weberpage{618}
%<<CONT
nicht Doppelwurzel ist. Denn die verschiedenen Wurzeln dieser Resolvente gehören zu conjugirten Gruppen, und wenn daher eine von diesen Wurzeln rational ist, so ist eine der conjugirten Gruppen metacyklisch. Wir haben also den Satz:

\gesperrt{V. Die nothwendige und hinreichende Bedingung für die Lösbarkeit der Gleichung $f(x)=0$ durch Radicale ist die, dass die Resolvente $\nu^{\mathrm{ten}}$ Grades $F(y)=0$, die man durch passende Wahl der Function $y$ so eingerichtet hat, dass sie keine Doppelwurzeln erhält, eine rationale Wurzel hat.}

Eine andere Form dieser Bedingung ergiebt sich auf folgende Weise.

Unter den Permutationen einer linearen Gruppe ist nur die identische, die irgend zwei Ziffern ungeändert lässt. Denn wenn $(z, az+b)$ zwei Ziffern nicht ändert, so muss die Congruenz $az+b \equiv z \pmod n$ zwei verschiedene Lösungen haben. Das ist aber nur möglich, wenn $a \equiv 1$, $b \equiv 0 \pmod n$ ist.

Ist also die Gruppe $P$ von $f(x)$ metacyklisch, so reducirt sie sich durch Adjunction von zwei beliebigen Wurzeln auf die Einheit, und folglich sind alle Wurzeln rational durch zwei beliebige unter ihnen ausdrückbar. Also:

\gesperrt{VI. Ist eine irreducible Gleichung von Primzahlgrad metacyklisch, so sind alle Wurzeln rational durch zwei beliebige unter ihnen ausdrückbar.}

Aber dieser Satz gilt auch umgekehrt, was wir folgendermaassen einfach beweisen können.

Es habe eine irreducible Gleichung vom Primzahlgrad $n$ die Eigenschaft, dass alle Wurzeln rational durch zwei beliebige unter ihnen ausdrückbar sind. Ist etwa $x_\nu = \psi(x_0, x_1)$ eine solche Darstellung, so können auf diese Relation alle Permutationen der Gruppe $P$ unserer Gleichung angewandt werden, und wenn also eine von diesen Permutationen $x_0$ und $x_1$ ungeändert lässt, so lässt sie auch alle $x_\nu$ ungeändert, d. h. es ist die identische Permutation. \gesperrt{Also enthält $P$ ausser der identischen Permutation keine, die zwei Ziffern nicht ändert.}

Enthält nun eine der Permutationen $\pi = cc_1 \ldots$ von $P$ zwei oder mehr verschiedene Cyklen $c, c_1 \ldots$ und ist $c$ vom Grade $h$, $c_1$ aber von einem höheren Grade, so ist $\pi^h = c_1 \ldots$, und $\pi^h$ lässt
%CONT>>
